% RESEARCH_DRAFT_SHA256: b7437a736dee43b59304d365b8dd62f31c24883ed00ee61db1043f4164e3f0f2
% Layout and source-reference wrappers expanded for the web renderer.
\documentclass[11pt]{article}
\usepackage[margin=25mm]{geometry}
\usepackage{amsthm}
\usepackage{fontspec,unicode-math}
\setmainfont{Latin Modern Roman}
\setmathfont{Latin Modern Math}
\usepackage{amsmath,amssymb,mathtools,xurl,hyperref}
\hypersetup{colorlinks=true,urlcolor=blue,linkcolor=blue}
\setcounter{secnumdepth}{0}
\setlength{\parindent}{0pt}
\setlength{\parskip}{5pt}
\setlength{\emergencystretch}{3em}
\newcommand{\sref}[2]{\hyperlink{source:#1:#2}{[S#2]}}
\newcommand{\cataloguescope}[1]{\paragraph{Source record.}#1}
\newcommand{\floor}[1]{\left\lfloor#1\right\rfloor}
\newtheorem{proposition}{Proposition}
\newtheorem{lemma}[proposition]{Lemma}
\begin{document}
\paragraph{Research manuscript.} OpenAI \texttt{gpt-6-astra}, reasoning \texttt{ultra}; 12 October 2026.
\paragraph{Preserved source body.} The complete catalogue mathematical text and sources follow with layout and source links expanded, including source qualifications and its historical preparation notes. New research begins at the section headed Research attempt.

\section{Reported Quantum Search for Binary Fusion Data}\label{statement}

\subsection{Definitions and mathematical statement}

\paragraph{Shared notation.}Standard number systems are $\mathbb N,\mathbb Z,\mathbb Q,\mathbb R,\mathbb C$. Specialized notation and index ranges follow the cited source and the definitions below.


A binary variable belongs to $\{0,1\}$. An efficient search problem additionally requires its constraint equations, input encoding, and success or completeness criterion.
\paragraph{Record scope and qualifications.}
The catalogue title is retained as metadata. Named certificates, contracts, bounds or frameworks in that title are not asserted as results of the original source extract.
The supplied registry flags this extraction as unrecoverable. It is displayed as a literal source fragment; the definitions provide context rather than a replacement formal statement.

The constraints involving the symbols Gamma are not supplied, so no definite binary-search target is inferred.

Source-record limitation: Source flags the statement as unrecoverable: The official report asks how to find binary fusion data efficiently and immediately presents the authors’ quantum algorithm as a solution; it notes that the algorithm produces examples rather than a mathematical classification. No separate clean statement is promoted; the original record remains publicly visible and retained for provenance.
The source text below is rendered literally where necessary; damaged formulas and composite extracts are not silently repaired into a new mathematical claim.
The supplied catalogue marks this record as unrecoverable. Its extract and terminology are documented, but a complete original problem has not been reconstructed.
\paragraph{Statement recorded in the supplied source.}
How to efficiently find Γkij \textbackslash{}in \textbackslash{}\{0, 1\textbackslash{}\} with solutions? The report passage is retained verbatim as the relevant context for this extracted problem. \href{https://doi.org/10.4171/owr/2025/34}{[S2]}
\paragraph{Source status.}
The supplied catalogue labels this source record open; that metadata was checked against the supplied snapshot on 2026-10-11. The label is preserved for provenance and is not an independent verification of present open status or mathematical correctness. The catalogue literature assessment is dated 2026-08-22.
\subsection{Short English statement}
The fragment asks for efficient search of binary fusion data.
\subsection{Sources}
\begin{itemize}
\item[S1] Ulam AI and the UnsolvedMath Contributors, UnsolvedMath: A Curated Collection of Open Mathematics Problems, catalogue entry. Catalogue curation and metadata: CC BY 4.0. The source extract is retained; definitions, English summaries and qualifications are editorial additions.
\url{https://huggingface.co/datasets/ulamai/UnsolvedMath}.
\item[S2] Subfactors and Applications (2026), Oberwolfach Reports, DOI 10.4171/owr/2025/34, pp. 1823--1886
\url{https://doi.org/10.4171/owr/2025/34}.
\item[S3] Oberwolfach Report 34/2025
\url{https://ems.press/content/serial-article-files/52239}.
\item Source-specific Oberwolfach Reports licences, including CC BY-SA 4.0 where applicable, are retained; see THIRD\_PARTY\_NOTICES.md.
\end{itemize}
\paragraph{Preparation and review.}
The original source and UnsolvedMath contributors retain attribution for the source extract. Codex (GPT-6; exact serving revision not exposed) assisted with definitions, English summaries, scope notes and formatting. No new human mathematical authorship or independently verified solution is asserted. Catalogue identity checked against the complete supplied registry and current repository. Distinct upstream records, including repeated questions and broad or damaged extracts, are retained for complete coverage of the source CS category. A repeated record is not a new independent conjecture. Complete mathematical-equivalence and current-literature review has not been performed. Preparation and final review follow contribution prompts C and D, with source gaps disclosed. Structural and rendering checks do not establish a new theorem or open status.
\label{end:statement}

\section{Research attempt}
\subsection{Literature, source context and scope}
The full original catalogue body and source qualifications are preserved above. The relevant primary workshop passages were read in the cached publisher report. Historical claims and subsequent attributed results are distinguished from the mathematical work below. A restricted example, standard lemma or known theorem is not counted as new progress on the entire original target. Missing assumptions and repeated or rhetorical extracts are not silently turned into different solved problems.

The original extract is not a standalone system of binary equations. The primary report explains fusion supports, forest restrictions, consistency equations and reflection positivity, then immediately describes a quantum example-generating algorithm. It expressly distinguishes example generation from mathematical classification. The catalogue's source-verification qualification is retained. This attempt tests the inference that evolution with negligible invalid-output probability necessarily explores most valid solutions; the report does not give enough Hamiltonian detail here to audit its claimed performance.
\subsection{Solution preservation does not imply coverage}
Let computational basis vectors correspond to finitely many binary strings. For a Hermitian matrix $H$, suppose the projector $P$ onto a subset of valid solutions commutes with $H$.
\begin{proposition}
A state initially supported in that subset stays there under $e^{itH}$ for every real $t$. Thus zero invalid-output probability alone does not ensure that other valid solutions can be found.
\end{proposition}
\begin{proof}
From $PH=HP$, induction gives $PH^j=H^jP$ for all nonnegative integers $j$. The finite-dimensional matrix exponential is the absolutely convergent power series $\sum_j(itH)^j/j!$, so $Pe^{itH}=e^{itH}P$. If $P\psi=\psi$, this implies $Pe^{itH}\psi=e^{itH}\psi$.

For an explicit binary-polynomial example, the constraint $uv=0$ has valid strings $00,01,10$. Let $H=|00\rangle\langle01|+|01\rangle\langle00|$, with zero action on $10,11$. Starting at $|00\rangle$ gives $\cos t\,|00\rangle+i\sin t\,|01\rangle$. Every measurement is valid, with zero error, but solution $10$ is never reached. This proves the asserted distinction without any asymptotic assumptions.
\end{proof}
\subsection{Even ideal uniform sampling has an output-coverage cost}
If independent samples were exactly uniform among $M\ge1$ solutions, the expected number required to see them all is $M H_M$, where $H_M=\sum_{j=1}^M1/j$. After $i$ distinct solutions have been seen, a new one appears with probability $(M-i)/M$ on each fresh draw; the waiting time has mean $M/(M-i)$ by summing its geometric tail. Summing these expectations for $i=0,\ldots,M-1$ proves the formula. For any fixed number $s$ of draws, the expected number unseen is $M(1-1/M)^s$ by linearity over the individual solutions. These are properties of the explicitly independent uniform model; correlated or biased Hamiltonian outputs need separate analysis.
\section{Verification and remaining gaps}
Matrix-power and exponential checks verify the three-solution example and exact occupancy formulas. This is an obstruction to a general coverage inference, not a counterexample to the source's particular unspecified Hamiltonian. To assess that algorithm one needs its implementable circuit/oracle costs, starting-solution access, reachability and sampling distribution, together with the full fusion constraints and positivity test. No new fusion data, efficient classification, or verified quantum speedup is claimed.

\section{References and source attribution}
Primary workshop report: \url{https://doi.org/10.4171/owr/2025/34}. Relevant passages read: OWR 2025/34, printed pages1855--1857, forest supports and reported quantum algorithm.

All retrievals were checked on 11 October 2026. Original question attribution remains in the preserved Sources section. External theorems are attributed and not claimed as new discoveries or fully verified merely by citation.


\paragraph{Integration review.} Codex (GPT-6; exact serving revision not exposed) reviewed the written arguments and expanded layout and source-reference wrappers. Repository metadata and the production web renderer were checked using Python 3.13 and JavaScript with MathJax 3.2.2. This model-assisted review is not independent human proof verification; no PDF compilation or formal proof checking is asserted.

\section{Completion Estimate}
\noindent\textbf{COMPLETION: 0\%}
No substantiated new progress toward the entire original target is claimed. The worked results and failed route are delimited above; a known result, illustrative model, or auxiliary obstruction is not counted as original-target completion.

\section{Final status and attribution}
\textbf{Attempt status: unresolved.} This describes this research attempt, not a new ruling on the source catalogue's status or on external literature claims.
The mathematics and this self-audit were generated by OpenAI \texttt{gpt-6-astra}, reasoning \texttt{ultra}, on 12 October 2026. Codex (GPT-6; exact serving revision not exposed) orchestrated the batch, read and reviewed the written mathematical arguments, independently reran the finite checks, and contributed formatting/repository review. No human mathematical authorship, independent human proof checking, or formal verification is claimed. The historical preparation paragraphs in the reproduced catalogue body concern that body only.

Contribution prompt A was applied to the complete definition and supplied source record, followed by the adversarial checks of prompt B. Local reproduction files are \texttt{batch026/checks.py}, \texttt{computations.json}, and the per-ID audit JSON. They were not public links at preparation time. Numerical checks are finite and do not prove unrestricted results. No LaTeX engine was found locally; compilation was not performed by this research agent.

\end{document}
